% A developed, fictional fourth-edition formatting demonstration.
% Build: xelatex example; biber example; xelatex example; xelatex example
% Or use `make sample` to keep auxiliary files under build/.
% TeX Gyre Termes is an explicit draft substitute. For submission, install
% licensed Times New Roman and change the option below to `submission`.
\documentclass[font=termes]{mbts}
\usepackage[american]{babel}
\usepackage{csquotes}
\usepackage[backend=biber,style=mbts]{biblatex}
\usepackage{mbts-bible}
\usepackage[hidelinks]{hyperref}
\addbibresource{example.bib}

\title{Learning to Listen:\\A Congregational Care Case Study}
\author{Morgan Reed}
\date{October 20, 2026}
\mbtssetup{course={HT 0000 Demonstration Seminar}}
\hypersetup{
  pdftitle={Learning to Listen: A Congregational Care Case Study},
  pdfauthor={Morgan Reed (fictional example)},
  pdfsubject={MBTS formatting demonstration; invented study and sources}
}
% This course-paper example includes its bibliography in the contents.
\AtBeginBibliography{\phantomsection
  \addcontentsline{toc}{section}{BIBLIOGRAPHY}}

\begin{document}
\maketitle

% These optional front-matter pages demonstrate a longer course paper.
% Keep only the pages your assignment requires.
\mbtsfrontmatter
\begin{mbtspage}*{About This Example}
This is a fictional research paper written to demonstrate the \texttt{mbtstex}
class. The student, congregation, interviews, statistics, secondary sources,
and source quotations are invented. The bibliography illustrates citation
formatting; it does not identify research you can consult. Scripture references
identify real biblical passages, but no Bible translation is quoted.

The example includes four heading levels, automatic footnotes, a block
quotation, Scripture references, numbered visuals, linked cross-references,
appendices, and an automatically sorted bibliography. The source files contain
comments at the main formatting decisions. The separate citation gallery
demonstrates additional source types that would make this paper considerably
longer and its fictional professor considerably less cheerful.

The typeface is TeX Gyre Termes, an explicit draft substitute. For submission,
use the class's \texttt{submission} option with licensed Times New Roman
installed, replace the invented material, and follow the assignment's
requirements. This explanatory page and the optional contents and visual lists
are included for demonstration. They are not universal course-paper requirements.
\end{mbtspage}
\tableofcontents
\listoftables
\listoffigures
\listofgraphs
\mbtsmainmatter

\section{Introduction}
Cedar Grove Church had a welcome team, a care team, and a committee appointed
to clarify the relationship between the two. It also had a coffee urn whose
operating instructions were better documented than its follow-up procedures.
This fictional case study asks how a small congregation might make listening
and follow-through more consistent without turning every conversation into a
new administrative department.

The argument is modest: a repeatable sequence of attentive conversation,
shared reflection, and an agreed next step can give volunteers a useful
structure for care. In the invented literature accompanying this example,
Anne Mercer distinguishes attentive presence from merely completing a
task.\autocite[24]{mercer-presence} Her distinction supplies the working
definition used here: a completed conversation creates room for the other
person's account before the volunteer proposes an action.\autocite[24]{mercer-presence}
The difference matters because a filled form can demonstrate that someone
held a pen. It cannot, by itself, demonstrate that someone listened.

The paper develops a theological frame, describes an illustrative training
process, and examines invented monthly records and interview excerpts.
The numbers show how to report a pattern with appropriate limits; they do not
establish that the proposed process would produce the same results elsewhere.
Appendix~\ref{app:protocol} supplies the interview prompts, and
appendix~\ref{app:records} supplies the full illustrative record.

% A deliberate page break keeps this major heading with its opening discussion.
\clearpage
\section{A Theological Frame for Attentive Care}
\subsection{Listening Before Answering}
James connects readiness to hear with restraint in speech and with action
\mbtsbible{James}{1:19--27}. For volunteers, the sequence poses a practical
question: has the speaker been understood before an answer has been prepared?

Mercer presents attention as an ordinary discipline, while Daniel Vale and
Miriam Bell emphasize an explicit next step.
\autocites[39--41]{mercer-presence}[62]{vale-bell}
These fictional sources frame listening and follow-through as one process.
An eloquent promise that no one remembers is still uncompleted, even when it
appears in beautifully formatted minutes.

\subsection{Neighbor Love and Particular Needs}
The encounter in \mbtsbibleref{Luke}{10:25--37} directs attention to a neighbor's
concrete situation. Cedar Grove's volunteers therefore ask what help is wanted
before recommending a familiar program. A meal, a ride, a later conversation,
and no immediate action are distinct possibilities.

Ruth Lin's fictional chapter places hospitality within the routines through
which a congregation receives people and remembers their concerns.
\autocite[87--89]{lin-hospitality} Here, hospitality names sustained attention
rather than an additional event. The calendar already has sufficient events
to qualify as a minor work of apocalyptic literature.

\section{Research Design}
\subsection{Setting and Participants}
The constructed setting consists of forty-eight households and a volunteer
team that grows from eight active members in January to twelve in December.
Six volunteers participate in the imagined interviews. They are selected to
illustrate different experiences of the process, not to represent a statistical
sample. No people were recruited, no interviews were conducted, and no
congregational records were collected for this example.

The design combines a monthly activity log with brief interviews about how
volunteers understand their work. A fictional dissertation by Leah Quill
provides the methodological model for reading the two kinds of material
together.\autocite[54--58]{quill-thesis} The log asks what was recorded; the
interview asks what the volunteer thought was happening. Agreement between
them may be informative, but a discrepancy can be equally useful.

\subsection{The Three-Step Practice}
Figure~\ref{fig:practice} presents the practice introduced in the hypothetical
training. A volunteer first listens without drafting a solution aloud, then
reflects back the concern, and finally agrees on an appropriate follow-up.
The arrows indicate a usual order rather than a prohibition on returning to
an earlier step. If a reflection is mistaken, the next move is to listen again.

% Native LaTeX drawing: no external image, graphics program, or TikZ needed.
\begin{mbtsvisual}{figure}{An Illustrative Care Conversation}
\label{fig:practice}
\centering
\setlength{\unitlength}{1pt}
\begin{picture}(420,46)
  \put(0,12){\framebox(110,26){Listen}}
  \put(110,25){\vector(1,0){45}}
  \put(155,12){\framebox(110,26){Reflect}}
  \put(265,25){\vector(1,0){45}}
  \put(310,12){\framebox(110,26){Follow Up}}
\end{picture}
\mbtsfiguresource{Author's invented training model.}
\mbtsfigurenote{The listener may return to an earlier step when clarification
is needed. The arrows do not represent measured effects.}
\end{mbtsvisual}

Practice sessions use ordinary situations: an interrupted conversation, a
request that needs clarification, and a promise that must be handed to another
volunteer. The team rehearses a brief summary and a specific next step. Ezra
Stone's fictional handbook describes a similar exercise in which volunteers
distinguish an invitation to help from an assumption that help is wanted.
\autocite[73]{stone-handbook}

\subsubsection{Recording a Conversation}
A record includes the month, whether a conversation was completed, and whether
an agreed follow-up was completed during that month. The example uses aggregate
counts, so a single household can appear in more than one conversation. These
records therefore cannot be interpreted as counts of distinct people served.

\paragraph{What Counts as Follow-Up}
A follow-up is counted only when the agreed action occurs. Scheduling a call
and making the call are separate events; the schedule alone does not complete
the action. The coding guide also distinguishes a newly recorded action from
an action carried over from a previous month.\autocite{cedar-guide}
This definition is intentionally less generous than the committee's initial
proposal to count sincere intentions, which would have produced exceptional
results at very little inconvenience.

\subsubsection{Reading the Interviews}
The invented interview excerpts are read for three themes: attention,
clarification, and follow-through. A response can receive more than one code.
The four authors of \textit{Listening in Community} provide the fictional
coding vocabulary used in this demonstration.\autocite[112--15]{park-team}
Appendix~\ref{app:protocol} includes prompts that invite examples and counterexamples,
including occasions when the process was unhelpful.

\section{Illustrative Findings}
\subsection{The Activity Record}
Table~\ref{tab:checkpoints} selects three checkpoints from the full record in
table~\ref{tab:monthly}. Recorded conversations rise from twelve in January
to twenty-seven in June, while the active volunteer count rises from eight to
ten. The table reports activity, not the quality of care or the number of
distinct households reached. Completed follow-ups may concern conversations
from an earlier month.

\begin{mbtsvisual}{table}{Selected Checkpoints in the Invented Record}
\label{tab:checkpoints}
\centering
\begin{mbtstabular}{lrrr}
  Month & Volunteers & Conversations & Follow-ups \\
  \hline
  January & 8 & 12 & 7 \\
  April   & 9 & 18 & 11 \\
  June    & 10 & 27 & 18 \\
  \hline
\end{mbtstabular}
\mbtsfiguresource{Invented data supplied in appendix~\ref{app:records}.}
\mbtsfigurenote{Volunteers are active team members; the other columns count
completed activities during the month.}
\end{mbtsvisual}

Graph~\ref{graph:activity} shows selected months across the full year. Its
zero baseline and common scale make the numerical increase visible without
presenting a small change as a cliff face. December's thirty-six conversations
are three times January's twelve. This description is arithmetic, not a claim
that the training caused a threefold improvement in ministry.

% A monochrome bar graph drawn with rules, all on the same scale.
\begin{mbtsvisual}{graph}{Recorded Conversations in Selected Months}
\label{graph:activity}
\centering
\renewcommand{\arraystretch}{1.25}
\begin{tabular}{l@{\quad}|@{}lr}
  & \multicolumn{1}{l}{Conversations (zero at left)} & Count \\
  January   & \rule{1in}{7pt}       & 12 \\
  March     & \rule{1.5in}{7pt}     & 18 \\
  May       & \rule{1.91667in}{7pt} & 23 \\
  July      & \rule{2.33333in}{7pt} & 28 \\
  September & \rule{2.58333in}{7pt} & 31 \\
  December  & \rule{3in}{7pt}       & 36 \\
\end{tabular}
\mbtsfiguresource{Invented monthly counts in table~\ref{tab:monthly}.}
\mbtsfigurenote{Each inch represents twelve conversations. Counts include
repeat conversations with the same household.}
\end{mbtsvisual}

The intervening months matter because a comparison of only two endpoints
can conceal a plateau or a reversal. April's conversation count, for example,
matches March's. Nothing in the log explains why. A responsible interpretation
leaves that question open rather than assigning the plateau to a spiritual
condition that the table was never designed to measure. Noah Wren makes the
same distinction between recorded activity and inferred meaning in the
fictional article used here.\autocite[49--51]{wren-measures}

% Start the interview discussion on a fresh page to keep the quotation intact.
\clearpage
\subsection{A Volunteer Describes the Change}
The following invented interview excerpt illustrates a shift from rapid
problem solving toward clarification. Volunteer June Avery describes an
early practice session:

\begin{mbtsquote}
I arrived with three suggestions and a clipboard. Before the other person
finished the first sentence, I had mentally assigned two volunteers and a
casserole. When I tried repeating back what I had heard, I discovered that
the request was for someone to sit and listen. The clipboard had prepared me
for almost everything except that. At the next meeting I wrote down the
person's own description before writing an action. We still arranged help,
but this time the help answered the question that had actually been asked.
The casserole remained available for emergencies.\autocite{avery-interview}
\end{mbtsquote}

Avery's account gives the proposed practice a concrete shape: reflection
interrupts an assumption and changes the next step. It does not establish
that the same change occurred in every conversation. Another invented
response describes the steps as awkward at first, especially when a
volunteer repeated a summary too mechanically. The case therefore treats
the sequence as a guide that requires judgment, not a script to be recited
until the other person recognizes the closing prayer.

\section{Interpretation and Limits}
\subsection{What the Pattern Can Support}
Within the constructed case, the records and interview excerpt are compatible
with a team becoming more deliberate about care. The numeric record supplies
a visible pattern, while the excerpt supplies one possible account of a
changed habit. Mercer argues in a second fictional book that practical
coordination can support a community's stated commitments when responsibility
is clear.\autocite[91--94]{mercer-meetings} That argument helps frame a useful
interpretation of the case without treating an activity count as a complete
description of faithfulness.

The strongest practical proposal is therefore to make the next step visible
and reviewable. At a later meeting, a volunteer can ask whether an agreed
action occurred and whether it met the expressed need. The team does not
need a committee to investigate whether it needs a committee. It needs a
shared understanding of who will do what, and an opportunity to correct a
mistake.\footnote{The case distinguishes coordination from the number of meetings
held. Compare \cite[97]{mercer-meetings}; \cite[66--68]{vale-bell}. A meeting can
produce a useful decision, but its duration is not a unit of pastoral care.}

\subsection{What the Pattern Cannot Support}
Several alternative explanations remain available even within the fictional
design. More active volunteers can produce more recorded conversations.
Volunteers may become more consistent in recording work they already did.
Households may make different requests at different times of year. Without
additional evidence, the table cannot separate those possibilities from a
change associated with training.\autocite[56]{wren-measures}

The six interview responses also cannot establish what all households
experienced. A volunteer's account of listening well may differ from the
other person's account of being heard. A fuller hypothetical study would
include those perspectives and actively examine cases that challenge its
preferred explanation. Quill's fictional dissertation provides an example of
this use of discrepant accounts.\autocite[121--24]{quill-thesis}
In this demonstration, restraint is part of the argument rather than a
paragraph added after the enthusiasm has spent the evidence.

\clearpage
\section{Conclusion}
This fictional case proposes a small, repeatable practice: listen, reflect,
and agree on a next step. The invented record illustrates how a writer can
describe activity, connect it to qualitative material, and preserve the
distinction between a pattern and an explanation. The theological frame
keeps attention on particular neighbors; the procedural frame makes a
promise easier to remember and review.

For Cedar Grove, the useful next question would be whether people receiving
care recognize the same improvement that volunteers describe. Until that
question is examined, a rising graph remains an invitation to inquire.
The coffee urn may continue to run on a fixed schedule. Congregational
care requires a little more attention.

\mbtsappendix{Illustrative Interview Protocol}
\label{app:protocol}
These prompts belong to the invented case study. They demonstrate an appendix
that supports the discussion without interrupting it. A researcher adapting
the example would develop questions appropriate to the actual study.

\begin{enumerate}
  \setlength{\itemsep}{0pt}
  \setlength{\parsep}{0pt}
  \item Describe a conversation in which you changed your understanding of
    what the other person wanted.
  \item What did you say or do to check that understanding?
  \item Describe an agreed follow-up. How did you know it was complete?
  \item When did the three-step practice feel awkward or unhelpful?
  \item Describe a conversation that does not fit your overall assessment
    of the training.
  \item What would the other person say about the same encounter?
\end{enumerate}

The imagined interviewer asks for examples before asking for an overall
evaluation. The order helps keep a general impression connected to a
particular event. Follow-up prompts seek clarification rather than praise
for the training. Neither coffee consumption nor the ability to pronounce
\textit{hermeneutics} on the first attempt is treated as an outcome measure.

\mbtsappendix{Complete Illustrative Activity Record}
\label{app:records}
Table~\ref{tab:monthly} contains all twelve invented monthly observations.
Its size makes an appendix the appropriate place for it. Follow-ups can
relate to earlier conversations; neither activity column counts unique people.

% This visual exceeds half a text page, so it belongs in an appendix.
\begin{mbtsvisual}{table}{Monthly Activity in the Fictional Congregation}
\label{tab:monthly}
\centering
\renewcommand{\arraystretch}{1.8}
\begin{mbtstabular}{lrrr}
  Month & Volunteers & Conversations & Follow-ups \\
  \hline
  January   & 8  & 12 & 7  \\
  February  & 8  & 14 & 8  \\
  March     & 9  & 18 & 10 \\
  April     & 9  & 18 & 11 \\
  May       & 10 & 23 & 15 \\
  June      & 10 & 27 & 18 \\
  July      & 11 & 28 & 19 \\
  August    & 11 & 30 & 21 \\
  September & 11 & 31 & 23 \\
  October   & 12 & 32 & 24 \\
  November  & 12 & 34 & 26 \\
  December  & 12 & 36 & 28 \\
  \hline
\end{mbtstabular}
\mbtsfiguresource{Author's invented dataset; no actual congregation was studied.}
\mbtsfigurenote{Volunteers are active team members, not new recruits.
Conversation and follow-up counts reset each month.}
\end{mbtsvisual}

\mbtsprintbibliography
\end{document}
